\section{Compresión}
\label{sec:compresion}

La compresión es el proceso de reducir el tamaño de los datos para disminuir el espacio de almacenamiento requerido y el tiempo de transmisión de los mismos. Se trata de una técnica fundamental en el procesamiento digital, con aplicaciones como la compresión de imágenes (por ejemplo, \texttt{JPEG}, \texttt{PNG}) y de audio (por ejemplo, \texttt{MP3}, \texttt{FLAC}).

En función del tratamiento de la información, existen dos tipos principales de compresión \cite{Kavitha2016A}:

\begin{itemize}
    \item Compresión sin pérdida (\textit{lossless}): permite reconstruir los datos originales de forma exacta a partir de los datos comprimidos, sin pérdida de información. Un ejemplo es el formato de audio \texttt{FLAC}, que comprime archivos de audio sin sacrificar calidad \cite{flac_spec}.
    \item Compresión con pérdida (\textit{lossy}): sacrifica parte de la información considerada menos perceptible o relevante para lograr una tasa de compresión mayor. En este caso, se busca que la diferencia entre el dato original y el reconstruido resulte imperceptible desde un punto de vista perceptual. Un ejemplo es el formato de audio \texttt{MP3}, que elimina componentes de alta frecuencia menos audibles para el oído humano para reducir el tamaño del archivo \cite{brandenburg1999mp3}.
\end{itemize}

Este último enfoque, el de compresión con pérdida, resulta especialmente relevante en la compresión de redes neuronales, donde se sacrifica precisión numérica en favor de la eficiencia computacional.

En el caso de la compresión de audio o de imágenes, el factor perceptual juega un papel crucial, ya que el objetivo es reducir el tamaño del archivo sin afectar significativamente la calidad percibida por el usuario.

En el contexto de la compresión de redes neuronales, el objetivo es reducir la cantidad de memoria y el costo computacional necesarios para almacenar y ejecutar el modelo, sin afectar significativamente las métricas de desempeño desarrolladas en la sección \ref{sec:metricas}.
